% Section 8.3, from lectures/L08/appendix/b_exercises.md.

\section{Exercises}
\label{c8:sec:exercises}\label{c8:app:b}

\begin{aside}
\textbf{How to work these.} These exercises supply the half \chapref{7} left out: a transport that
can be scripted, so that the provided host suite can drive the real \code{Uart} with no hardware,
plus the blocking helpers on top of the non-blocking core. You write the stub and the helpers. A
small demo, \file{main.cpp}, is provided too, listed in full in Exercise~\ref{c8:ex:2.2}.

\textbf{How to check your work.} The provided host suite, \file{fw/test/uart/uart_test.cpp}, run by
\code{make test}, is the check, and it is the first one the C++ half of the book has been able to
run. It is guarded on all four driver headers --- \file{register_map.hpp} and \file{uart.hpp} from
Chapters~6 and~7, \file{stub.hpp} and \file{blocking.hpp} from here --- so it reports there is
nothing to do until the last of them exists; adding the two files below is what switches it on.

Every exercise in this chapter is code, so Appendix~\ref{sol:ch} has nothing to add here.
\end{aside}

Work in \file{fw/}. You will add:

\begin{console}
include/
    driver/
        uart/
            blocking.hpp
        transport/
            stub.hpp
source/
    main.cpp            (provided; the host demo, copied from the last exercise)
test/
    uart/
        uart_test.cpp   (provided; the host suite, already in the repository)
\end{console}

The stub joins the transport interface in \code{driver::transport} (the demo may use a small
\code{app} namespace). These conventions apply to \file{include/} and \file{source/}, which have to
cross-compile for the ATmega; \file{test/} and \file{include/arch/test/} are host-only and never
reach avr-gcc, which is why the provided suites next to them use \code{<cstdint>}, \code{std::},
\code{namespace driver::uart::test} and \code{[[nodiscard]]} freely. Follow the same AVR-portable
conventions as in \chapref{6} (\secref{c6:sec:exercises}): \code{<stdint.h>} and bare \code{uint8_t}
and \code{size_t} (no \code{std::}), nested namespace blocks (not
\code{namespace driver::uart { ... }}), and no \code{[[nodiscard]]}. Build the application with
\code{make build} and the tests with \code{make test}.

% ------------------------------------------------------------------------------------------
\exerciseset{The Scripted Stub}
\label{c8:set:1}

\exercise{\code{driver::transport::Stub}}{C++}
\label{c8:ex:1.1}
In \file{include/driver/transport/stub.hpp}, build a header-only \code{driver::transport::Stub} that
implements \code{driver::transport::Interface} so that the driver can be tested with no hardware. It
is the wire, scripted, and it has three jobs.

The provided test suite binds to this stub by name, so the names below are part of the exercise
rather than suggestions: a stub with the right behaviour but different member names will not compile
against \file{test/uart/uart_test.cpp}. It needs a default constructor, since every test builds one
with \code{transport::Stub stub{};}.

That suite opens with four \code{TransportStub} cases that check the stub itself, before any driver
case runs. They are there because every \code{Uart} case reads what the stub was scripted to reply,
so a stub bug surfaces as a \emph{driver} failure: queue the four bytes least significant first, for
instance, and the suite reports \code{Uart.WriteWhenReadyPushesTxData} red, pointing at
\code{readReg()}, which is correct. Get the \code{TransportStub} cases green first and the failures
below are genuinely yours.

It must \textbf{record} every call. Append each byte passed to \code{transfer()} to a fixed-size
\code{uint8_t} buffer with a running length, where a capacity of 100 bytes is plenty for these
tests. Expose that record through \code{txLen()}, returning the number of bytes captured, and
\code{txByte(index)}, returning the byte at a position and \code{0} for an index past the end. Both
must be \textbf{\code{const}}, because the tests inspect the stub through a
\code{const transport::Stub&}. Keep two \code{uint16_t} counters as well, one bumped by
\code{begin()} and one by \code{end()}, exposed as \code{beginCalls()} and \code{endCalls()}, also
\code{const}. These are not optional: the suite asserts on \code{beginCalls()} in eight separate cases
and on \code{endCalls()} in two, because they are what proves a transaction was framed exactly once and that the two stay balanced.

It must \textbf{play back} bytes on \code{transfer()}, returning the next byte from a preloaded
response buffer tracked by an index, and returning \code{0x00} once that buffer is exhausted. This is
the \code{MISO} the driver reads.

And it must offer helpers to script that response. \code{injectRxByte(uint8_t)} queues a single
byte, and \code{injectRxWord(uint32_t)} queues a whole register value \textbf{most significant
first}, so that a \code{readReg} sees exactly that value. \code{clearRxData()} resets the reply
buffer and its index, so that one test can script several reads.

Give the two a \emph{different name} rather than overloading one name on width. An integer literal
converts equally well to \code{uint8_t} and to \code{uint32_t}, so an overload pair would make
\code{stub.injectRx(0x41)} ambiguous and force every call site to spell out the type of its argument.
Two names cost nothing and keep the call sites readable, which is why the UART stub of \chapref{6} is
built the same way.

Two fixed \code{uint8_t} buffers with a length each, one for what the driver sent and one for the
scripted replies, are all it takes; no dynamic containers are needed. Both are \textbf{linear
buffers, not FIFOs}: the record buffer is a write-only append log the test reads back afterwards, and
the reply buffer is read front to back through an advancing index. No circular wraparound, no
pop-with-shift; just reset each length, and the read index, to zero at the start of a test.

Keep it dependency-injected: a test constructs a \code{Stub}, constructs a \code{Uart} over it,
drives the driver, and then asserts against the \code{Stub}'s record.

% ------------------------------------------------------------------------------------------
\exerciseset{Using the Driver}
\label{c8:set:2}

\exercise{Blocking helpers}{C++}
\label{c8:ex:2.1}
The interface's \code{write()} and \code{read()} are non-blocking by design. Build the thin blocking
convenience on top, in a new header \file{include/driver/uart/blocking.hpp}, as free functions in the
\code{driver::uart} namespace. Keep the same AVR-portable style as the rest of the driver
(\code{<stdint.h>} and bare types, nested namespace blocks). Both take \code{Interface&}, not
\code{Uart&}, so they work over any implementation; \code{uart.write} and \code{uart.read} dispatch
to the concrete driver at run time.

\task{a) Write a byte, blocking}
Add a free function \code{writeBlocking()} that sends one byte, waiting until the transmitter accepts
it. It takes a \code{driver::uart::Interface&}, the driver, and a \code{uint8_t}, the byte to send;
it spins on \code{uart.write(byte)} until that returns \code{true} and then returns nothing. Mark it
\code{inline}, since it is defined in a header and that avoids a one-definition-rule violation once
more than one file includes it, and mark it \code{noexcept}.

\task{b) Read a byte, blocking}
Add a free function \code{readBlocking()} that receives one byte, waiting until one is available. It
takes a \code{driver::uart::Interface&}, the driver, and a reference to a \code{uint8_t} where the
received byte is stored; it spins on \code{uart.read(byte)} until that returns \code{true} and then
returns nothing. It is \code{inline} and \code{noexcept} for the same reasons.

Each is a one-line spin loop, so header-only \code{inline} functions are all you need; no
\file{.cpp}. They belong above the driver, not inside it: the core stays deterministic and testable,
and the spinning lives where a caller opts into it.

\exercise{The provided demo}{C++}
\label{c8:ex:2.2}
Use this test program to exercise communication end to end. It is \textbf{provided}, like the host
suite; you do not write it. Create \file{source/main.cpp} and copy it in as it stands:

\begin{cppcode}
/**
 * @brief Application entry point.
 */
#include <stdint.h>

#include "driver/transport/stub.hpp"
#include "driver/uart/blocking.hpp"
#include "driver/uart/register_map.hpp"
#include "driver/uart/uart.hpp"

using namespace driver;

/**
 * @brief Application entry point.
 *
 *        On the host, the demo uses a scripted transport::Stub because no real
 *        UART is available. It reports TX ready and returns one received byte so
 *        the blocking calls can complete. On the ATmega328P (L09), the real
 *        transport reports the actual hardware status.
 *
 *        Each register read transfers one command byte followed by four data
 *        bytes, so scripted responses begin with a command-phase placeholder.
 *        Because writes also consume scripted bytes, each response is installed
 *        immediately before use and clearRxData() resets the script between
 *        phases.
 *
 * @return 0 on termination of the program.
 */
int main()
{
    constexpr uint16_t baudDiv{27U};
    constexpr uint8_t txByte{0x7FU};
    constexpr uint8_t dummyCmd{0U};

    constexpr uint32_t rxData{txByte};
    constexpr uint32_t txReady{static_cast<uint32_t>(1U << uart::status::TX_READY)};
    constexpr uint32_t rxValid{static_cast<uint32_t>(1U << uart::status::RX_VALID)};

    transport::Stub transport{};
    uart::Uart uart{transport};
    uart.configure(baudDiv);

    // Report TX ready, then send a byte.
    transport.injectRxByte(dummyCmd);
    transport.injectRxWord(txReady);
    uart::writeBlocking(uart, txByte);

    // Report RX valid and hand back the byte, then receive it.
    transport.clearRxData();
    transport.injectRxByte(dummyCmd);
    transport.injectRxWord(rxValid);
    transport.injectRxByte(dummyCmd);
    transport.injectRxWord(rxData);

    uint8_t rxByte{};
    uart::readBlocking(uart, rxByte);
    return 0;
}
\end{cppcode}

It shows the API end to end: it constructs a \code{Uart} over a \code{driver::transport::Stub},
\code{configure()}s it, then uses your \code{writeBlocking} and \code{readBlocking} to send a byte and
receive one back. Because the blocking calls spin until the transport reports ready, the demo
scripts the stub just before each call so that it completes instead of spinning forever. On the
target there is nothing to script: \chapref{9} writes a separate freestanding \code{main} in
\file{fw/avr/} that builds the same \code{Uart} over \code{AvrSpi}, and the real transport reports
the actual hardware status.

Read it, then run \code{make build} and \code{./uart_firmware}: it exits cleanly once your
\code{Uart}, \code{Stub} and blocking helpers are correct. The demo is illustration, not a test; the
real checking is \code{make test}.

\subsection*{What's ahead}
\chapref{9} replaces the stub with \code{AvrSpi}, the real SPI on the ATmega328P, and the driver, the
interface and the register map all run unchanged over the wire. The suite you just turned green runs
unchanged too, which is the return on putting the seam where \chapref{6} put it.
